\begin{figure}[htbp]
  \centering
  \begin{tikzpicture}
    \begin{axis}[
      xbar,
      width=0.78\textwidth,
      height=4.5cm,
      bar width=18pt,
      xlabel={Стоимость, руб. (логарифмическая шкала)},
      xmode=log,
      xmin=10, xmax=500000,
      log ticks with fixed point,
      ytick=data,
      symbolic y coords={{ИИ-генерация курса},{Подрядная разработка}},
      y tick label style={font=\small},
      nodes near coords,
      nodes near coords align={horizontal},
      every node near coord/.append style={font=\footnotesize},
      xmajorgrids=true,
      grid style={dashed, gray!40},
      enlarge y limits=0.5,
    ]
    \addplot[fill=teal!60, draw=teal!80] coordinates {
      (75,{ИИ-генерация курса})
    };
    \addplot[fill=red!40, draw=red!60] coordinates {
      (150000,{Подрядная разработка})
    };
      \node[font=\small, text=teal!60!black, anchor=west]
        at (rel axis cs:0.56,0.78) {$\times$2000 разница};
    \end{axis}
  \end{tikzpicture}
  \caption{Unit-экономика ИИ-методиста: себестоимость vs рыночная стоимость}
  \label{fig:unit-economics}
\end{figure}
